\section{Discussion}

This study supports three practical conclusions.

First, for Katharevousa parsing in this data regime, multilingual transformer encoders remain the most reliable path for structural syntax quality. Second, LLM-assisted data generation can be scientifically productive when operational safeguards are treated as first-class methodology rather than implementation details. Third, framework comparisons (mBERT, Stanza pretrained/custom) are useful not only for ranking but for diagnosing model-family behavior under historical-language constraints.

From a digital humanities perspective, the contribution is not just model accuracy. The pipeline operationalizes a reproducible way to move from difficult archival text to reusable language technology while preserving a transparent record of methodological choices and failure modes.
